% ============================================================
%  Manual de Capacitación — Plataforma KikiBrows
%  Compilar con: pdflatex (Overleaf compatible)
% ============================================================

\documentclass[a4paper,11pt]{article}

\usepackage[utf8]{inputenc}
\usepackage[T1]{fontenc}
\usepackage[spanish]{babel}
\usepackage{geometry}
\geometry{a4paper, top=25mm, bottom=28mm, left=25mm, right=25mm}
\usepackage{lmodern}
\usepackage{microtype}
\usepackage{parskip}
\setlength{\parindent}{0pt}
\setlength{\parskip}{6pt}

\usepackage{booktabs}
\usepackage{longtable}
\usepackage{array}
\usepackage{tabularx}
\newcolumntype{L}[1]{>{\raggedright\arraybackslash}p{#1}}
\newcolumntype{C}[1]{>{\centering\arraybackslash}p{#1}}

\usepackage{xcolor}
\definecolor{mainblack}{RGB}{17,17,17}
\definecolor{gold}{RGB}{184,150,62}
\definecolor{darkgray}{RGB}{44,44,44}
\definecolor{midgray}{RGB}{85,85,85}
\definecolor{bordergray}{RGB}{210,210,210}
\definecolor{softbg}{RGB}{247,247,245}
\definecolor{goldbg}{RGB}{245,237,217}
\definecolor{warnbg}{RGB}{255,244,236}
\definecolor{warnorange}{RGB}{192,90,0}
\definecolor{checkgreen}{RGB}{30,110,30}
\definecolor{navpublic}{RGB}{220,235,220}
\definecolor{navalumna}{RGB}{220,228,245}
\definecolor{navadmin}{RGB}{245,237,217}
\definecolor{navroot}{RGB}{30,30,30}

\usepackage{titlesec}
\titleformat{\section}
  {\large\bfseries\color{mainblack}}
  {\color{gold}\thesection.}{0.5em}{}
  [\vspace{-4pt}{\color{gold}\rule{\linewidth}{1.5pt}}\vspace{2pt}]
\titleformat{\subsection}
  {\normalsize\bfseries\color{darkgray}}{\thesubsection}{0.5em}{}
\titleformat{\subsubsection}
  {\small\bfseries\color{midgray}}{}{0em}{}
\titlespacing*{\section}{0pt}{18pt}{6pt}
\titlespacing*{\subsection}{0pt}{12pt}{4pt}
\titlespacing*{\subsubsection}{0pt}{8pt}{2pt}

\usepackage{fancyhdr}
\pagestyle{fancy}
\fancyhf{}
\fancyhead[L]{\small\color{midgray}Manual de Capacitaci\'on --- Plataforma KikiBrows}
\fancyhead[R]{\small\color{midgray}RAWSEC SOLUTIONS SpA}
\fancyfoot[C]{\small\color{midgray}\thepage}
\renewcommand{\headrulewidth}{0.3pt}
\renewcommand{\headrule}{\color{bordergray}\hrule width\headwidth height\headrulewidth}

\usepackage{mdframed}
\newmdenv[
  backgroundcolor=goldbg, linecolor=gold, linewidth=3pt,
  topline=false, bottomline=false, rightline=false,
  innertopmargin=10pt, innerbottommargin=10pt,
  innerleftmargin=14pt, innerrightmargin=12pt,
  skipabove=10pt, skipbelow=8pt,
]{infobox}
\newmdenv[
  backgroundcolor=warnbg, linecolor=warnorange, linewidth=3pt,
  topline=false, bottomline=false, rightline=false,
  innertopmargin=10pt, innerbottommargin=10pt,
  innerleftmargin=14pt, innerrightmargin=12pt,
  skipabove=10pt, skipbelow=8pt,
]{warnbox}
\newmdenv[
  backgroundcolor=navpublic, linecolor=checkgreen, linewidth=3pt,
  topline=false, bottomline=false, rightline=false,
  innertopmargin=8pt, innerbottommargin=8pt,
  innerleftmargin=14pt, innerrightmargin=12pt,
  skipabove=8pt, skipbelow=6pt,
]{alumnabox}
\newmdenv[
  backgroundcolor=navadmin, linecolor=gold, linewidth=3pt,
  topline=false, bottomline=false, rightline=false,
  innertopmargin=8pt, innerbottommargin=8pt,
  innerleftmargin=14pt, innerrightmargin=12pt,
  skipabove=8pt, skipbelow=6pt,
]{adminbox}

\usepackage{enumitem}
\setlist[enumerate,1]{label=\textbf{\arabic*.}, leftmargin=1.8em, itemsep=4pt}
\setlist[itemize,1]{leftmargin=1.6em, itemsep=3pt}

% Árbol de navegación (folder tree)
\usepackage{forest}

\usepackage{hyperref}
\hypersetup{
  colorlinks=true, linkcolor=gold, urlcolor=gold,
  pdftitle={Manual de Capacitacion -- KikiBrows},
  pdfauthor={RAWSEC SOLUTIONS SpA},
}

% ============================================================
\begin{document}

% ══════════════════════════════════════════════════════════
%  PORTADA
% ══════════════════════════════════════════════════════════
\thispagestyle{empty}
\vspace*{18mm}

\begin{center}
  {\small\bfseries\color{gold}MANUAL DE CAPACITACI\'ON}\\[14pt]
  {\LARGE\bfseries\color{mainblack}
    Plataforma de Cursos Profesionales\\[6pt]KikiBrows}\\[12pt]
  {\large\color{midgray}Gu\'ia de uso para administradoras y alumnas}\\[28pt]
  {\color{bordergray}\rule{0.45\linewidth}{0.4pt}}\\[22pt]
\end{center}

\begin{center}
\renewcommand{\arraystretch}{1.4}
\begin{tabular}{@{}ll@{}}
  \textcolor{midgray}{\small Preparado por}   & \textbf{RAWSEC SOLUTIONS SpA} \\
  \textcolor{midgray}{\small Para}             & \textbf{KikiBrows SpA} \\[4pt]
  \textcolor{midgray}{\small Fecha}            & Junio de 2026 \\
  \textcolor{midgray}{\small Versi\'on}        & 1.0 \\
\end{tabular}
\end{center}

\vfill
\begin{center}
  {\color{bordergray}\rule{\linewidth}{0.3pt}}\\[5pt]
  {\small\color{midgray}
    Documento de uso interno --- KikiBrows SpA --- Junio 2026}
\end{center}
\newpage

% ══════════════════════════════════════════════════════════
%  TABLA DE CONTENIDOS
% ══════════════════════════════════════════════════════════
\tableofcontents
\newpage

% ══════════════════════════════════════════════════════════
%  1. INTRODUCCIÓN
% ══════════════════════════════════════════════════════════
\section{Introducci\'on}

Este manual describe c\'omo usar la plataforma de cursos KikiBrows, tanto desde el
rol de \textbf{administradora} como desde el rol de \textbf{alumna}. La plataforma
permite vender y gestionar cursos de est\'etica profesional, hacer seguimiento del
progreso de las alumnas, recibir pagos autom\'aticos y coordinar consultas en vivo.

\subsection*{Roles del sistema}

\begin{center}
\renewcommand{\arraystretch}{1.4}
\begin{tabular}{L{3cm} L{11.5cm}}
\toprule
\textbf{Rol} & \textbf{Capacidades} \\
\midrule
\textbf{Administradora} &
  Crea y edita cursos, gestiona usuarias, revisa transacciones,
  crea slots de consultas en vivo y califica entregas. \\
\textbf{Alumna} &
  Compra cursos, accede al contenido, realiza quizzes, env\'ia tareas,
  agenda consultas en vivo y descarga certificados. \\
\bottomrule
\end{tabular}
\end{center}

\begin{infobox}
\textbf{Acceso:} La plataforma funciona desde cualquier navegador web moderno
(Chrome, Safari, Edge, Firefox). No requiere instalar ning\'un programa.
La URL de acceso es la direcci\'on web del sitio de KikiBrows.
\end{infobox}

% ══════════════════════════════════════════════════════════
%  2. ÁRBOL DE NAVEGACIÓN
% ══════════════════════════════════════════════════════════
\section{Mapa de Navegaci\'on}

El siguiente diagrama muestra el flujo real de navegaci\'on de la plataforma:
qu\'e p\'aginas existen, c\'omo se llega a cada una y qu\'e rol puede acceder.
El punto de entrada es siempre la p\'agina de \textbf{Inicio} (p\'ublica);
desde ah\'i se accede al cat\'alogo, a la compra de cursos y al inicio de sesi\'on.

\begin{infobox}
\textbf{\textcolor{mainblack}{¿C\'omo se entra a cada zona?}}
Existen dos formas de llegar a cada zona privada, ambas determinadas por el
\textbf{rol} de la cuenta:
\begin{itemize}[leftmargin=1.6em, itemsep=3pt, topsep=2pt]
  \item \textbf{Al iniciar sesi\'on} en \emph{login}, el sistema deriva
        autom\'aticamente: la \textbf{administradora} llega directo a su
        \emph{Panel de Administraci\'on}, y la \textbf{alumna} vuelve a la
        p\'agina de Inicio.
  \item \textbf{Con la sesi\'on ya abierta}, desde el men\'u de usuaria del Inicio
        se usa la opci\'on \textbf{<<Mi Cuenta>>}, que redirige seg\'un el rol:
        a la \textbf{alumna} la lleva a su cuenta (\emph{account}) y a la
        \textbf{administradora} la lleva a su \emph{Panel de Administraci\'on}.
\end{itemize}
Las p\'aginas de cada zona est\'an protegidas: si una alumna intenta abrir una
ruta de administraci\'on, el sistema la devuelve al Inicio; si alguien sin sesi\'on
intenta abrir una p\'agina privada, se le env\'ia a \emph{login}.
\end{infobox}

\newpage
{\large\bfseries\color{mainblack}Mapa de Navegaci\'on --- Plataforma KikiBrows}

\vspace{4pt}
{\small\color{midgray}
  \textcolor{checkgreen!70!black}{\rule{7mm}{3pt}}~Acceso p\'ublico\quad
  \textcolor{blue!50!black}{\rule{7mm}{3pt}}~Portal Alumna\quad
  \textcolor{gold}{\rule{7mm}{3pt}}~Panel Administradora
}

\vspace{10pt}

\begin{center}
\footnotesize
\begin{forest}
  for tree={
    grow'=0,
    s sep=3pt,
    draw,
    rounded corners=2pt,
    inner xsep=4pt, inner ysep=3pt,
    edge={bordergray, line width=0.8pt},
    anchor=west,
    child anchor=west,
    parent anchor=south,
    calign=first,
    edge path={
      \noexpand\path [draw, \forestoption{edge}]
        (!u.south west) +(8pt,0) |- (.child anchor)\forestoption{edge label};
    },
    before typesetting nodes={
      if n=1{insert before={[,phantom]}}{}
    },
    fit=band,
    before computing xy={l=16pt},
  }
  [\textbf{Plataforma KikiBrows}, fill=navroot, text=white, draw=navroot
    [\textbf{ACCESO P\'UBLICO} \textcolor{checkgreen!60!black}{\scriptsize(sin sesi\'on)},
        fill=navpublic, draw=checkgreen!70!black,
        for descendants={fill=navpublic!45, draw=checkgreen!55!black}
      [Inicio / Landing \textcolor{midgray}{\scriptsize(index.html)}
        [Cat\'alogo de cursos]
        [Vista previa de curso \textcolor{midgray}{\scriptsize(course-preview.html)}
          [Comprar $\rightarrow$ Pasarela Banchile $\rightarrow$ Confirmaci\'on de Pago
              \textcolor{midgray}{\scriptsize(payment-confirmation.html)}]
        ]
        [Men\'u de usuaria \textcolor{midgray}{\scriptsize(visible con sesi\'on iniciada)}
          [<<Mi Cuenta>> $\rightarrow$ redirige seg\'un el rol:
            [Alumna $\rightarrow$ Mi Cuenta \textcolor{midgray}{\scriptsize(account.html)}]
            [Administradora $\rightarrow$ Panel de Administraci\'on \textcolor{midgray}{\scriptsize(adminPanel.html)}]
          ]
          [Mis Cursos $\rightarrow$ Portal Alumna \textcolor{midgray}{\scriptsize(cursosAlumn.html)}]
          [Consultas en Vivo \textcolor{midgray}{\scriptsize(consultasAlumn.html)}]
        ]
      ]
      [Pol\'itica de Privacidad \textcolor{midgray}{\scriptsize(politica-privacidad.html)}]
      [Iniciar sesi\'on \textcolor{midgray}{\scriptsize(login.html)}
        [Registrarse \textcolor{midgray}{\scriptsize(registrate.html)}]
        [Recuperar contrase\~na \textcolor{midgray}{\scriptsize(recover.html $\rightarrow$ resetPassword.html)}]
        [Redirecci\'on autom\'atica seg\'un el rol
          [Alumna $\rightarrow$ Inicio (men\'u de alumna)]
          [Administradora $\rightarrow$ Panel de Administraci\'on]
        ]
      ]
    ]
    [\textbf{PORTAL ALUMNA} \textcolor{blue!50!black}{\scriptsize(requiere sesi\'on)},
        fill=navalumna, draw=blue!50!black,
        for descendants={fill=navalumna!45, draw=blue!45!black}
      [Mis Cursos \textcolor{midgray}{\scriptsize(cursosAlumn.html)}
        [Aula Virtual / Clase \textcolor{midgray}{\scriptsize(claseAlumn.html)}
          [Contenido: Video {\textperiodcentered} Texto {\textperiodcentered} PDF]
          [Quiz]
          [Tarea / Entrega pr\'actica]
        ]
      ]
      [Mi Cuenta \textcolor{midgray}{\scriptsize(account.html)}
        [Historial de Compras \textcolor{midgray}{\scriptsize(historialCompras.html)}]
      ]
      [Mis Certificados \textcolor{midgray}{\scriptsize(historialCertificados.html)}]
      [Consultas en Vivo \textcolor{midgray}{\scriptsize(consultasAlumn.html)}]
    ]
    [\textbf{PANEL ADMINISTRADORA} \textcolor{gold}{\scriptsize(requiere rol admin)},
        fill=navadmin, draw=gold!80!black,
        for descendants={fill=navadmin!45, draw=gold!70!black}
      [Dashboard \textcolor{midgray}{\scriptsize(adminPanel.html)}]
      [Usuarios \textcolor{midgray}{\scriptsize(usersGest.html)}]
      [Gesti\'on de Cursos \textcolor{midgray}{\scriptsize(gestionCursos.html)}
        [Crear / Editar Curso \textcolor{midgray}{\scriptsize(creaCurso.html)}
          [Gestor de M\'odulos y Clases \textcolor{midgray}{\scriptsize(gestorModulos.html)}
            [Previsualizar Curso \textcolor{midgray}{\scriptsize(previsualizaCurso.html)}]
          ]
        ]
      ]
      [Transacciones \textcolor{midgray}{\scriptsize(adminTransa.html)}]
      [Calendario / Consultas en vivo \textcolor{midgray}{\scriptsize(adminCalendar.html)}]
      [Revisiones y Feedback \textcolor{midgray}{\scriptsize(revYFeedback.html)}]
      [Cambiar Contrase\~na \textcolor{midgray}{\scriptsize(adminProfilePassword.html)}]
      [Volver al Sitio \textcolor{midgray}{\scriptsize($\rightarrow$ index.html)}]
    ]
  ]
\end{forest}
\end{center}

% ══════════════════════════════════════════════════════════
%  3. ACCESO AL SISTEMA
% ══════════════════════════════════════════════════════════
\section{Acceso al Sistema}

\subsection{Iniciar sesi\'on}

\begin{enumerate}
  \item Ingresa a la URL del sitio KikiBrows.
  \item Haz clic en \textbf{Iniciar sesi\'on}.
  \item Escribe tu correo electr\'onico y haz clic en \textbf{Continuar}.
  \item Revisa tu bandeja de entrada: recibir\'as un \textbf{c\'odigo de verificaci\'on
        de 6 d\'igitos}.
  \item Ingresa el c\'odigo en la pantalla. Tienes 10 minutos antes de que expire.
  \item La plataforma te redirige autom\'aticamente seg\'un tu rol
        (Panel de administraci\'on o Mis Cursos).
\end{enumerate}

\begin{infobox}
\textbf{Sin contrase\~na:} El sistema usa verificaci\'on por correo electr\'onico
(c\'odigo OTP) en vez de contrase\~na fija. Si el c\'odigo no llega, revisa la
carpeta de spam o espera un minuto y sol\'icita uno nuevo.
\end{infobox}

\subsection{Registro de nueva alumna}

\begin{enumerate}
  \item En la p\'agina de inicio o login, haz clic en \textbf{Registrarse}.
  \item Completa el formulario:
        nombre, apellido, correo electr\'onico.
  \item Lee y acepta la \textbf{Pol\'itica de Privacidad} (obligatorio).
  \item Haz clic en \textbf{Crear cuenta}. Llegar\'a un c\'odigo de verificaci\'on
        al correo ingresado.
  \item Ingresa el c\'odigo para activar la cuenta.
\end{enumerate}

\subsection{Recuperar acceso}

\begin{enumerate}
  \item En la p\'agina de login, haz clic en \textbf{Olvidé mi acceso} o
        \textbf{Recuperar}.
  \item Ingresa tu correo electr\'onico y env\'ia el formulario.
  \item Recibir\'as un enlace de restablecimiento en tu correo.
  \item Haz clic en el enlace e ingresa una nueva contrase\~na.
\end{enumerate}

% ══════════════════════════════════════════════════════════
%  4. PANEL DE LA ADMINISTRADORA
% ══════════════════════════════════════════════════════════
\section{Panel de la Administradora}

Al iniciar sesi\'on con una cuenta administradora, la plataforma redirige al
\textbf{Panel de Administraci\'on}. Desde aqu\'i se accede a todas las funciones
de gesti\'on.

\begin{adminbox}
\textbf{Acceso exclusivo:} Las p\'aginas de administraci\'on solo son accesibles
para usuarias con rol \emph{Admin}. Si una alumna intenta acceder a estas rutas,
el sistema la redirige al portal de alumnas.
\end{adminbox}

% ─────────────────────────────────────────
\subsection{Gesti\'on de Cursos}

\subsubsection{Crear un nuevo curso}

\begin{enumerate}
  \item Ve a \textbf{Gesti\'on de Cursos} → \textbf{Crear Curso}.
  \item Completa los campos:
  \begin{itemize}
    \item \textbf{Nombre del curso} (obligatorio)
    \item \textbf{Descripci\'on}
    \item \textbf{Precio} (en pesos chilenos)
    \item \textbf{D\'ias de acceso}: cu\'antos d\'ias tendr\'a acceso la alumna
          despu\'es de comprar (valor por defecto: 180 d\'ias)
    \item \textbf{Estado}: \emph{Borrador} (solo visible para admin),
          \emph{Publicado} (visible para alumnas), \emph{Pausado}
  \end{itemize}
  \item Sube la \textbf{imagen de portada} del curso
        (formatos PNG, JPG o WEBP; m\'aximo 2\,MB).
  \item Si quieres que el curso aparezca en el carrusel de la p\'agina de inicio,
        activa la opci\'on \textbf{Mostrar en carrusel} e indica la posici\'on.
  \item Haz clic en \textbf{Guardar Curso}.
  \item Una vez guardado, aparece la secci\'on de \textbf{M\'odulos}. Agrega
        al menos un m\'odulo escribiendo su nombre y haciendo clic en \textbf{Agregar}.
\end{enumerate}

\begin{warnbox}
\textbf{Antes de agregar clases} debes guardar el curso. El bot\'on para editar
clases dentro de un m\'odulo solo se activa despu\'es del primer guardado.
\end{warnbox}

\subsubsection{Editar y reordenar m\'odulos}

\begin{itemize}
  \item Para cambiar el nombre de un m\'odulo: haz clic en el \'icono de edici\'on
        junto al nombre, modif\'icalo y guarda.
  \item Para cambiar el orden: arrastra el m\'odulo con el \'icono de reordenamiento
        (\emph{drag \& drop}) a la posici\'on deseada.
  \item Para eliminar un m\'odulo: haz clic en el \'icono de papelera.
        \textbf{Advertencia:} esto elimina el m\'odulo y todas sus clases.
  \item Para editar las clases de un m\'odulo: haz clic en \textbf{Gestionar Clases}.
\end{itemize}

% ─────────────────────────────────────────
\subsection{Gestor de M\'odulos y Clases}

Desde la pantalla del \textbf{Gestor de M\'odulos}, se agregan y editan las clases
de cada m\'odulo.

\subsubsection{Crear una nueva clase}

\begin{enumerate}
  \item Haz clic en \textbf{Agregar Clase}.
  \item Elige el \textbf{tipo de clase}:
\end{enumerate}

\renewcommand{\arraystretch}{1.3}
\begin{longtable}{L{2.8cm} L{11.7cm}}
\toprule
\textbf{Tipo} & \textbf{Descripci\'on y campos} \\
\midrule
\endhead
\bottomrule
\endfoot

\textbf{Video} &
  Sube un archivo de video (MP4, WEBM, MOV, HEVC o MKV; m\'aximo 50\,MB).
  Ingresa nombre, descripci\'on, duraci\'on en minutos.
  El sistema muestra una barra de progreso durante la subida. \\

\textbf{Texto} &
  Escribe el contenido directamente en el editor de texto enriquecido
  (admite negritas, listas, t\'itulos). Sin archivos adjuntos. \\

\textbf{PDF} &
  Sube un archivo PDF (m\'aximo 50\,MB). Las alumnas pueden verlo
  en pantalla pero no descargarlo directamente. \\

\textbf{Quiz} &
  Configura las instrucciones, el porcentaje de aprobaci\'on
  (por defecto 70\,\%) y agrega las preguntas con sus opciones
  de respuesta, marcando la opci\'on correcta. \\

\textbf{Pr\'actica} &
  Las alumnas deben subir un archivo (video o imagen) como entrega.
  Configura las instrucciones, puntaje total, nota de aprobaci\'on,
  y opcionalmente un video demostrativo e im\'agenes de referencia. \\

\end{longtable}

\begin{enumerate}
  \setcounter{enumi}{2}
  \item Completa los campos seg\'un el tipo de clase.
  \item Haz clic en \textbf{Guardar Clase}.
  \item Para reordenar clases dentro del m\'odulo, arr\'astralas con el \'icono
        de reordenamiento.
\end{enumerate}

\begin{infobox}
\textbf{Duraci\'on del curso:} La duraci\'on total del curso (suma de todas las
clases) se calcula autom\'aticamente y se muestra en el editor.
\end{infobox}

% ─────────────────────────────────────────
\subsection{Gesti\'on de Usuarias}

Desde \textbf{Gesti\'on de Usuarias} se administra el acceso de todas las personas
registradas en la plataforma.

\subsubsection{Ver y buscar usuarias}

\begin{itemize}
  \item La tabla muestra nombre, correo, rol, estado y cursos activos.
  \item Usa el \textbf{buscador} para filtrar por nombre o correo.
  \item Usa el \textbf{filtro de rol} para ver solo alumnas o solo administradoras.
\end{itemize}

\subsubsection{Crear nueva usuaria}

\begin{enumerate}
  \item Haz clic en \textbf{Nueva Usuaria}.
  \item Completa nombre, apellido, correo y contrase\~na
        (m\'inimo 6 caracteres).
  \item Selecciona el rol: \emph{Alumna} o \emph{Admin}.
  \item Opcionalmente asigna un curso desde el men\'u desplegable.
  \item Haz clic en \textbf{Crear}.
\end{enumerate}

\subsubsection{Editar, bloquear o eliminar usuaria}

\begin{itemize}
  \item \textbf{Editar:} haz clic en el \'icono de edici\'on de la fila,
        modifica los datos y guarda. Desde aqu\'i tambi\'en puedes
        asignar un curso manualmente.
  \item \textbf{Bloquear/Desbloquear:} activa o desactiva el interruptor de
        bloqueo. Una usuaria bloqueada no puede iniciar sesi\'on.
  \item \textbf{Eliminar:} haz clic en el \'icono de papelera. El sistema
        pide confirmaci\'on. \textbf{Esta acci\'on es irreversible} y borra
        la usuaria junto con sus cursos y pagos registrados.
\end{itemize}

\begin{warnbox}
\textbf{No puedes editarte a ti misma.} El interruptor de bloqueo est\'a
desactivado en tu propia fila para evitar bloquearte accidentalmente.
\end{warnbox}

% ─────────────────────────────────────────
\subsection{Transacciones}

Desde \textbf{Transacciones} puedes ver todos los pagos realizados en la plataforma.

\begin{itemize}
  \item La tabla muestra: folio, alumna, curso, monto, m\'etodo de pago,
        estado (\emph{Pagado}, \emph{Pendiente}, \emph{Fallido}) y fecha.
  \item Usa los \textbf{filtros} para buscar por estado, m\'etodo o rango de fechas.
  \item Los pagos realizados con \textbf{Banchile Pagos} incluyen el c\'odigo
        de autorizaci\'on del banco.
\end{itemize}

\begin{infobox}
\textbf{Pagos autom\'aticos:} Cuando una alumna paga un curso con Banchile,
el sistema genera la inscripci\'on autom\'aticamente y env\'ia un comprobante
por correo. No se requiere acci\'on manual.
\end{infobox}

% ─────────────────────────────────────────
\subsection{Consultas en Vivo}

Desde \textbf{Consultas en Vivo} (panel admin) creas los horarios disponibles
para que las alumnas reserven sus consultas. El sistema crea la reuni\'on
de Zoom autom\'aticamente.

\subsubsection{Crear un nuevo slot de consulta}

\begin{enumerate}
  \item Haz clic en \textbf{Nuevo Slot}.
  \item Selecciona la \textbf{fecha} y los \textbf{horarios} de inicio y fin.
  \item Indica el n\'umero m\'aximo de \textbf{cupos}.
  \item Haz clic en \textbf{Crear}. El sistema genera autom\'aticamente
        la reuni\'on en Zoom y guarda los enlaces.
\end{enumerate}

\subsubsection{Ver detalles de un slot}

\begin{itemize}
  \item Cada slot muestra el n\'umero de cupos ocupados vs.\ disponibles
        y el estado: \emph{Disponible}, \emph{Lleno}, \emph{Completado}
        o \emph{Cancelado}.
  \item Haz clic en un slot para ver la lista de alumnas inscritas con
        nombre, correo y estado de confirmaci\'on.
  \item Desde el detalle puedes copiar el \textbf{link de host} (para ti)
        y el \textbf{link de alumna} (para compartir) con un clic.
\end{itemize}

\begin{warnbox}
\textbf{Antes de eliminar un slot:} si hay alumnas inscritas, el sistema muestra
una advertencia con el n\'umero de inscritas. Conf\'irmalo solo si notificaste
previamente a las alumnas afectadas.
\end{warnbox}

% ─────────────────────────────────────────
\subsection{Revisi\'on y Feedback}

Desde \textbf{Revisi\'on y Feedback} revisas las tareas pr\'acticas enviadas
por las alumnas.

\begin{enumerate}
  \item Selecciona la entrega que deseas revisar.
  \item Visualiza el archivo enviado por la alumna (video o imagen).
  \item Ingresa tu \textbf{calificaci\'on} (n\'umero) y el \textbf{feedback}
        (comentario escrito).
  \item Selecciona el estado: \textbf{Aprobada} o \textbf{Rechazada}.
  \item Haz clic en \textbf{Guardar}. La alumna ver\'a el resultado y tu
        comentario la pr\'oxima vez que abra esa clase.
\end{enumerate}

\begin{infobox}
Si la entrega es \textbf{Rechazada}, la alumna puede volver a subir un nuevo
intento. El historial de todos los intentos queda guardado.
\end{infobox}

% ─────────────────────────────────────────
\subsection{Mi Perfil}

Desde \textbf{Mi Perfil} puedes actualizar tu contrase\~na de acceso. Ingresa
la contrase\~na actual y la nueva, conf\'irmala y guarda.

% ══════════════════════════════════════════════════════════
%  5. PORTAL DE LA ALUMNA
% ══════════════════════════════════════════════════════════
\section{Portal de la Alumna}

\begin{alumnabox}
\textbf{Acceso:} Despu\'es de iniciar sesi\'on, la alumna llega directamente
a \textbf{Mis Cursos}, su p\'agina principal donde ve todos los cursos
que ha comprado o que le han sido asignados.
\end{alumnabox}

% ─────────────────────────────────────────
\subsection{Mis Cursos}

La p\'agina \textbf{Mis Cursos} muestra una tarjeta por cada curso al que
la alumna tiene acceso, con:

\begin{itemize}
  \item Imagen de portada y nombre del curso.
  \item \textbf{Barra de progreso} con el porcentaje completado.
  \item \textbf{Tiempo restante de acceso} (en d\'ias). Si quedan menos de
        7 d\'ias, el indicador se muestra en rojo.
\end{itemize}

\renewcommand{\arraystretch}{1.3}
\begin{longtable}{L{3.8cm} L{10.7cm}}
\toprule
\textbf{Estado del bot\'on} & \textbf{Situaci\'on} \\
\midrule
\endhead
\bottomrule
\endfoot
Comenzar Curso      & La alumna a\'un no ha abierto ninguna clase (0\,\%). \\
Continuar Aprendiendo & La alumna est\'a en progreso (entre 1\,\% y 99\,\%). \\
Revisar Contenido   & El curso est\'a completado (100\,\%). \\
Acceso Expirado     & La fecha de acceso venci\'o; el bot\'on est\'a bloqueado. \\
\end{longtable}

\begin{itemize}
  \item Usa el \textbf{buscador} para filtrar cursos por nombre.
  \item Haz clic en la tarjeta o en el bot\'on para entrar al curso.
\end{itemize}

% ─────────────────────────────────────────
\subsection{Dentro de un Curso}

Al entrar a un curso, la pantalla se divide en:

\begin{itemize}
  \item \textbf{Panel lateral izquierdo:} lista de m\'odulos y clases con
        iconos de estado (completada, en progreso, bloqueada).
  \item \textbf{\'Area principal:} contenido de la clase seleccionada.
  \item \textbf{Barra de progreso:} porcentaje global del curso, visible
        en la parte superior.
\end{itemize}

Para navegar usa los botones \textbf{Anterior} y \textbf{Siguiente} o haz
clic directamente en una clase del panel lateral.

\begin{warnbox}
\textbf{Clases bloqueadas:} algunas clases requieren completar la anterior
para desbloquearse. Un \'icono de candado indica que a\'un no est\'an
disponibles.
\end{warnbox}

\subsubsection{Ver un video}

\begin{enumerate}
  \item Haz clic en una clase de tipo \textbf{Video} en el panel lateral.
  \item El video se reproduce en el \'area principal. Puedes pausar, avanzar
        y ajustar el volumen.
  \item La plataforma recuerda en qu\'e segundo qued\'o el video
        si cierras y vuelves m\'as tarde.
  \item La clase se marca como \textbf{completada} autom\'aticamente cuando
        ves al menos el 90\,\% del video.
\end{enumerate}

\subsubsection{Leer un texto o PDF}

\begin{enumerate}
  \item Selecciona la clase de tipo \textbf{Texto} o \textbf{PDF}.
  \item Lee el contenido en pantalla.
  \item Haz clic en \textbf{Marcar como completada} cuando termines.
\end{enumerate}

\subsubsection{Realizar un quiz}

\begin{enumerate}
  \item Selecciona la clase de tipo \textbf{Quiz}.
  \item Lee las instrucciones y responde cada pregunta.
  \item Haz clic en \textbf{Enviar respuestas}.
  \item Ver\'as tu calificaci\'on y, para cada pregunta, cu\'al era
        la respuesta correcta.
  \item Si no apruebas, puedes intentarlo nuevamente.
        El historial de intentos queda visible.
\end{enumerate}

\begin{infobox}
El porcentaje m\'inimo para aprobar un quiz est\'a definido por la administradora
al crear la clase (generalmente 70\,\%). Aprobar el quiz desbloquea la
siguiente clase si est\'a configurado como requisito.
\end{infobox}

\subsubsection{Enviar una tarea (Pr\'actica)}

\begin{enumerate}
  \item Selecciona la clase de tipo \textbf{Pr\'actica}.
  \item Lee las instrucciones y revisa las im\'agenes o video de referencia
        si los hay.
  \item Haz clic en \textbf{Seleccionar archivo} y elige tu video (MP4, WEBM, MOV;
        m\'ax.\ 50\,MB) o imagen (PNG, JPG, WEBP; m\'ax.\ 2\,MB).
  \item Opcionalmente agrega un comentario para la instructora.
  \item Haz clic en \textbf{Enviar entrega}. Se mostrar\'a una barra de progreso
        durante la subida.
  \item La entrega queda en estado \textbf{Pendiente de Revisi\'on} hasta que
        la instructora la califique.
\end{enumerate}

\renewcommand{\arraystretch}{1.3}
\begin{longtable}{L{3.8cm} L{10.7cm}}
\toprule
\textbf{Estado de entrega} & \textbf{Significado} \\
\midrule
\endhead
\bottomrule
\endfoot
Pendiente de Revisi\'on & La entrega fue enviada y est\'a esperando revisi\'on. \\
Entrega Aprobada        & La instructora aprob\'o la tarea. Se muestra el feedback. \\
Entrega Rechazada       & La instructora rechaz\'o la tarea con comentario. Puedes reenviar. \\
\end{longtable}

% ─────────────────────────────────────────
\subsection{Consultas en Vivo}

\begin{enumerate}
  \item Ve a \textbf{Consultas en Vivo} desde el men\'u principal.
  \item La p\'agina muestra los horarios disponibles con fecha, hora y
        n\'umero de cupos libres.
  \item Haz clic en \textbf{Reservar} en el slot que desees.
  \item Selecciona el curso sobre el que quieres consultar
        (o elige \emph{Consulta General}).
  \item Confirma la reserva.
  \item En la secci\'on \textbf{Mis Reservas} ver\'as tu slot con el
        bot\'on \textbf{Ver link clase} cuando el enlace de Zoom est\'e disponible.
\end{enumerate}

\begin{infobox}
Los cupos son limitados. Si el slot aparece como \emph{Lleno}, monitorea la
p\'agina ya que los cupos pueden liberarse si otra alumna cancela.
\end{infobox}

% ─────────────────────────────────────────
\subsection{Historial de Compras y Certificados}

\subsubsection{Historial de Compras}

\begin{enumerate}
  \item Ve a \textbf{Historial de Compras} desde el men\'u.
  \item Ver\'as la lista de todos tus pagos con: folio, curso, monto, fecha
        y estado (\emph{Pagado}, \emph{Pendiente}, \emph{Rechazado}).
  \item Haz clic en \textbf{Descargar Comprobante} para obtener el PDF
        del comprobante de pago.
\end{enumerate}

\subsubsection{Mis Certificados}

\begin{enumerate}
  \item Ve a \textbf{Mis Certificados} desde el men\'u.
  \item Los certificados aparecen cuando has completado el 100\,\% del curso.
  \item Haz clic en \textbf{Descargar} para obtener tu certificado en PDF.
  \item Cada certificado tiene un \textbf{c\'odigo de verificaci\'on \'unico}
        que permite a terceros validar su autenticidad.
\end{enumerate}

% ══════════════════════════════════════════════════════════
%  6. COMPRA DE UN CURSO (FLUJO DE PAGO)
% ══════════════════════════════════════════════════════════
\section{Compra de un Curso}

\begin{enumerate}
  \item En la p\'agina de inicio, haz clic en el curso que quieres comprar.
  \item Ver\'as la \textbf{vista previa del curso}: descripci\'on, contenido
        por m\'odulos y precio.
  \item Haz clic en \textbf{Comprar}.
  \item Si no tienes sesi\'on iniciada, el sistema te pedir\'a que
        inicies sesi\'on o te registres primero.
  \item Ser\'as redirigida al portal de pago de \textbf{Banchile Pagos}.
  \item Completa los datos de tu tarjeta y confirma el pago.
  \item Al finalizar, la plataforma te redirige a la pantalla de
        \textbf{Confirmaci\'on de Pago}. Si el pago fue aprobado:
  \begin{itemize}
    \item El curso aparece en \textbf{Mis Cursos} de inmediato.
    \item Recibes un \textbf{comprobante de pago} por correo electr\'onico.
  \end{itemize}
\end{enumerate}

\begin{warnbox}
\textbf{Si el pago fue rechazado:} la plataforma muestra el motivo del rechazo
(tarjeta inv\'alida, fondos insuficientes, etc.). Puedes intentarlo nuevamente
con otra tarjeta desde la vista previa del curso.
\end{warnbox}

% ══════════════════════════════════════════════════════════
%  7. PREGUNTAS FRECUENTES
% ══════════════════════════════════════════════════════════
\section{Preguntas Frecuentes}

\renewcommand{\arraystretch}{1.5}
\begin{longtable}{L{6cm} L{8.5cm}}
\toprule
\textbf{Pregunta} & \textbf{Respuesta} \\
\midrule
\endhead
\bottomrule
\endfoot

No me llega el c\'odigo de verificaci\'on. &
  Revisa la carpeta de spam. Si no aparece, espera 2 minutos y
  sol\'icita un nuevo c\'odigo desde la misma pantalla. \\

La alumna no puede acceder a su curso. &
  Verifica en \emph{Gesti\'on de Usuarias} que la cuenta no est\'e
  bloqueada y que la inscripci\'on no haya expirado. \\

Una clase aparece bloqueada para la alumna. &
  Significa que la clase anterior es requisito y a\'un no fue completada.
  Revisa si hay un quiz pendiente de aprobar. \\

\textquestiondown C\'omo s\'e que la alumna complet\'o el curso? &
  En \emph{Gesti\'on de Usuarias} la fila de la alumna muestra su
  c\'odigo de certificado. Tambi\'en puedes ver el estado en sus cursos. \\

Elimin\'e una clase por error. &
  La eliminaci\'on de clases es irreversible. Deber\'as recrear la clase
  y volver a subir el contenido. \\

El video tard\'o mucho en subir. &
  Los videos se suben directamente al servidor. El tiempo depende del
  tama\~no del archivo y la velocidad de internet.
  Los videos de hasta 50\,MB pueden tardar varios minutos. \\

\textquestiondown D\'onde est\'an los pagos de una alumna espec\'ifica? &
  Ve a \emph{Transacciones} y busca por nombre o correo en el filtro. \\

La alumna dice que no le lleg\'o el comprobante de pago. &
  Verifica en \emph{Transacciones} que el estado sea \emph{Pagado}.
  El comprobante se genera autom\'aticamente; revisa spam. La alumna tambi\'en
  puede descargarlo desde su \emph{Historial de Compras}. \\

\end{longtable}

% ══════════════════════════════════════════════════════════
%  9. MANTENIMIENTO DE LA PLATAFORMA
% ══════════════════════════════════════════════════════════
\newpage
\section{Mantenimiento de la Plataforma}

La plataforma KikiBrows funciona sobre servicios en la nube (Supabase, Cloudflare,
Resend, Zoom API) que operan de manera aut\'onoma una vez configurados. Sin embargo,
como toda plataforma en producci\'on, requiere tareas peri\'odicas de revisi\'on y
mantenci\'on para garantizar su correcto funcionamiento.

Este cap\'itulo distingue entre las tareas que puede realizar la administradora
directamente desde el panel web, y las tareas que requieren acceso t\'ecnico a
los servicios de infraestructura.

% ─────────────────────────────────────────────────────────
\subsection{Tareas mensuales a cargo de la administradora}

Las siguientes actividades no requieren conocimiento t\'ecnico y pueden ejecutarse
desde el \textbf{Panel de Administraci\'on} de la plataforma:

\subsubsection*{Revisi\'on de transacciones}
\begin{itemize}
  \item Verificar en \emph{Transacciones} que los pagos del mes tienen estado
        \emph{Pagado}. Los pagos en estado \emph{Pendiente} por m\'as de 48\,h
        pueden indicar un problema con la pasarela.
  \item Exportar o anotar los montos para control contable mensual.
\end{itemize}

\subsubsection*{Revisi\'on de avances y certificaciones}
\begin{itemize}
  \item Verificar que las alumnas que completaron cursos hayan recibido su certificado.
        En \emph{Gesti\'on de Usuarias} se muestra el c\'odigo de certificado por fila.
  \item Revisar en \emph{Revisiones y Feedback} si hay entregas pendientes de calificar.
\end{itemize}

\subsubsection*{Gesti\'on de consultas en vivo}
\begin{itemize}
  \item Publicar los \emph{slots} de consulta del mes siguiente en
        \emph{Calendario / Consultas} antes del inicio de cada mes.
  \item Verificar que las reservas agendadas tienen el enlace Zoom correspondiente.
\end{itemize}

\subsubsection*{Revisi\'on de usuarias}
\begin{itemize}
  \item Detectar usuarias inactivas por per\'iodo prolongado (sin inicio de sesi\'on)
        y decidir si deben ser bloqueadas o contactadas.
  \item Si se cambia la contrase\~na de acceso al correo institucional o a Zoom,
        notificar a RAWSEC SOLUTIONS SpA para actualizar las credenciales en el sistema.
\end{itemize}

% ─────────────────────────────────────────────────────────
\subsection{L\'imites de almacenamiento a monitorear}

La plataforma utiliza el plan gratuito de Supabase Storage. Los l\'imites por tipo
de archivo son los siguientes:

\begin{center}
\renewcommand{\arraystretch}{1.4}
\begin{tabular}{@{}L{5cm}C{3.5cm}L{6cm}@{}}
  \toprule
  \textbf{Tipo de archivo} & \textbf{L\'imite por archivo} & \textbf{D\'onde aplica} \\
  \midrule
  Videos de clases         & 50\,MB & Todas las clases de tipo video \\
  Documentos PDF           & 20\,MB & Material descargable de clases \\
  Im\'agenes (portadas)    &  2\,MB & Portadas de cursos y m\'odulos \\
  Entregas de alumnas      & 50\,MB & Archivos subidos por alumnas en tareas \\
  \bottomrule
\end{tabular}
\end{center}

\begin{warnbox}
\textbf{Si un archivo supera el l\'imite,} el sistema rechazar\'a la subida con
un mensaje de error. Comprimir el video o reducir la resoluci\'on antes de volver
a intentarlo. Si el contenido del curso lo requiere, se puede solicitar a
RAWSEC SOLUTIONS SpA la migraci\'on a Supabase Pro (almacenamiento ampliado).
\end{warnbox}

El almacenamiento total del proyecto en el plan gratuito de Supabase es de
\textbf{500\,MB}. Cuando la plataforma crezca en n\'umero de cursos o alumnas,
ser\'a necesario migrar al plan pago de Supabase (actualmente \$25\,USD/mes),
lo que incluye almacenamiento ilimitado, mayor ancho de banda y funcionalidades
adicionales.

% ─────────────────────────────────────────────────────────
\subsection{Tareas t\'ecnicas que requieren soporte especializado}

Las siguientes tareas requieren acceso a los servicios de infraestructura
(Supabase, Cloudflare, Resend, Zoom API) y \textbf{no deben intentarse sin
conocimiento t\'ecnico}, ya que un error puede dejar la plataforma inoperativa.

\begin{itemize}
  \item \textbf{Actualizaci\'on de credenciales de pasarela de pago:} cualquier
        cambio en las claves de Banchile, Resend, o Zoom debe aplicarse en el
        panel de \emph{Secrets} de Supabase Edge Functions. Un valor incorrecto
        bloquea los pagos o el env\'io de correos.

  \item \textbf{Limpieza de archivos hu\'erfanos en Storage:} videos o im\'agenes
        eliminados desde el panel no siempre se borran autom\'aticamente del
        bucket de Supabase. Peridicamente deben eliminarse los archivos sin
        referencia activa.

  \item \textbf{Migraci\'on de plan Supabase:} si el proyecto supera los 500\,MB
        de almacenamiento o los l\'imites de peticiones del plan gratuito, debe
        actualizarse el plan desde el panel de Supabase con una tarjeta de
        cr\'edito asociada al proyecto.

  \item \textbf{Renovaci\'on del dominio en Cloudflare:} el dominio de la
        plataforma tiene vencimiento peri\'odico. Su no renovaci\'on hace
        inaccesible el sitio.

  \item \textbf{Revisi\'on de Edge Functions:} en caso de fallos en el env\'io
        de correos o procesamiento de pagos, los logs de las funciones
        (\texttt{send-receipt-email}, \texttt{send-consultation-email},
        \texttt{banchile-payment-*}) deben revisarse desde la consola de Supabase.
\end{itemize}

% ─────────────────────────────────────────────────────────
\subsection{Servicio de mantenci\'on con RAWSEC SOLUTIONS SpA}

\begin{infobox}
\textbf{Opci\'on de mantenci\'on mensual}

KikiBrows SpA puede contratar a \textbf{RAWSEC SOLUTIONS SpA} para que ejecute
las tareas t\'ecnicas descritas en la secci\'on anterior, incluyendo:

\begin{itemize}
  \item Monitoreo mensual de almacenamiento y uso de la plataforma.
  \item Limpieza de archivos hu\'erfanos en Storage.
  \item Actualizaci\'on de credenciales ante cambios de proveedores.
  \item Revisi\'on de logs de pagos y funciones de correo.
  \item Renovaci\'on del dominio web (Cloudflare) cuando corresponda.
  \item Soporte reactivo ante incidentes t\'ecnicos (dentro del mes contratado).
\end{itemize}

Para cotizar el servicio de mantenci\'on mensual, contactar a:\\
\textbf{RAWSEC SOLUTIONS SpA} --- \texttt{rawsecsolutions@gmail.com}
\end{infobox}

% ══════════════════════════════════════════════════════════
%  FIRMA Y CIERRE
% ══════════════════════════════════════════════════════════
\vspace{10mm}
\begin{center}
  {\color{bordergray}\rule{\linewidth}{0.3pt}}\\[6pt]
  {\small\color{midgray}
    Manual de Capacitaci\'on --- Plataforma KikiBrows \,\textperiodcentered\,
    RAWSEC SOLUTIONS SpA \,\textperiodcentered\, Junio 2026\\
    Para soporte t\'ecnico contactar a RAWSEC SOLUTIONS SpA.}
\end{center}

\end{document}
